\chapter{Cross-Cloud Service Equivalence}
\index{cross-cloud equivalence}\index{service mapping}\index{multi-cloud}%
This appendix is a living reference that maps conceptually equivalent services and
ideas across the seven platforms studied in this book series: Amazon Web Services
(AWS), Microsoft Azure, Microsoft Fabric, Google Cloud Platform (GCP), MongoDB,
Snowflake, and Databricks. The names differ, but the underlying data-engineering
concepts---durable object storage, tiered lifecycle economics, immutability, managed
relational engines, elastic compute, and disciplined data organization---recur
everywhere.

\begin{notebox}
This appendix grows with each chapter. Every new chapter contributes the equivalence
tables introduced in its cross-cloud callout, so by the end of the book you will hold
a single consolidated Rosetta Stone for moving fluently between clouds. Where a clean
one-to-one mapping does not exist, the prose notes the closest analog and the caveat.
\end{notebox}

\section{Object and Data-Lake Storage}
\index{object storage!equivalence}\index{storage classes!equivalence}%
The foundational layer covered in Chapter~1. Every platform offers durable, HTTP-addressable
object storage with tiered pricing, lifecycle transitions, versioning, and write-once-read-many
(WORM) controls.

\begin{center}
\footnotesize
\begin{tabular}{@{}p{2.8cm}p{3.2cm}p{3.2cm}p{3.2cm}@{}}
\toprule
\textbf{Capability} & \textbf{AWS} & \textbf{Azure} & \textbf{GCP} \\
\midrule
Object storage & Amazon S3 & ADLS Gen2 / Blob & Cloud Storage \\
Hot tier & S3 Standard & Hot & Standard \\
Infrequent access & S3 Standard-IA & Cool & Nearline \\
Archive (instant) & Glacier Instant Retrieval & Cold & Coldline \\
Deep archive & Glacier Deep Archive & Archive & Archive \\
Lifecycle rules & Lifecycle policies & Lifecycle management & Object Lifecycle Mgmt \\
Versioning & S3 Versioning & Blob versioning & Object Versioning \\
Immutability (WORM) & Object Lock & Immutable blob storage & Bucket Lock \\
Cross-region copy & CRR / SRR & Object replication & Dual/multi-region + Transfer \\
Encryption at rest & SSE-S3 / SSE-KMS & SSE + CMK (Key Vault) & Google-managed / CMEK \\
\bottomrule
\end{tabular}
\end{center}

\noindent The three remaining platforms approach storage differently.
\textbf{Microsoft Fabric} exposes a single logical lake, \emph{OneLake}, and uses
\emph{shortcuts} to virtualize S3, ADLS Gen2, and Google Cloud Storage in place without copying.
\textbf{Snowflake} abstracts storage entirely: table data lives in cloud object storage as
immutable, columnar \emph{micro-partitions} that the service manages on your behalf.
\textbf{MongoDB} persists documents through the \emph{WiredTiger} storage engine, balancing an
internal cache against the operating-system page cache rather than exposing storage tiers.

\section{Managed Relational Databases}
\index{relational databases!equivalence}%
The focus of Chapter~2 for the three hyperscalers. Each provides a fully managed relational
engine, a cloud-native high-performance variant, high availability, read replicas, and a
migration service.

\begin{center}
\footnotesize
\begin{tabular}{@{}p{2.8cm}p{3.2cm}p{3.2cm}p{3.2cm}@{}}
\toprule
\textbf{Capability} & \textbf{AWS} & \textbf{Azure} & \textbf{GCP} \\
\midrule
Managed relational & Amazon RDS & Azure SQL Database & Cloud SQL \\
Cloud-native engine & Amazon Aurora & SQL Hyperscale & AlloyDB \\
Managed PostgreSQL & RDS / Aurora PostgreSQL & Azure DB for PostgreSQL & Cloud SQL / AlloyDB \\
High availability & Multi-AZ & Zone-redundant HA & Regional HA \\
Read scaling & Read Replicas & Read replicas / geo-replicas & Read Replicas \\
Global / geo & Aurora Global Database & Active geo-replication & Cross-region replicas \\
Serverless compute & Aurora Serverless v2 & Azure SQL Serverless & AlloyDB (elastic) \\
Migration service & DMS + SCT & Azure Database Migration Svc & Database Migration Svc \\
\bottomrule
\end{tabular}
\end{center}

\noindent \textbf{Snowflake}, \textbf{Fabric}, and \textbf{MongoDB} are not managed relational
engines, but they solve overlapping problems: Snowflake and Fabric provide SQL analytics over
warehouse/lakehouse storage, while MongoDB provides a distributed document store whose schema
design (Chapter~2) replaces rigid relational normalization with intentional, access-driven modeling.

\section{Elastic Compute and Analytical Warehouses}
\index{elastic compute!equivalence}\index{virtual warehouse!equivalence}%
Snowflake's Chapter~2 topic---separating compute from storage and scaling it elastically---has a
counterpart in every analytics platform.

\begin{center}
\footnotesize
\begin{tabular}{@{}p{2.8cm}p{3.2cm}p{3.2cm}p{3.2cm}@{}}
\toprule
\textbf{Concept} & \textbf{Snowflake} & \textbf{AWS} & \textbf{Azure / GCP} \\
\midrule
Compute unit & Virtual Warehouse & Redshift RA3 / Serverless & Synapse pool / BigQuery slots \\
Elastic concurrency & Multi-cluster warehouse & Concurrency scaling & Autoscale / on-demand \\
Pause when idle & Auto-suspend / resume & Pause / resume & Auto-pause / serverless \\
Billing unit & Credits & RPU / node-hours & DWU / slot-hours \\
Result reuse & Result cache & Result cache & Result / BI cache \\
\bottomrule
\end{tabular}
\end{center}

\noindent \textbf{Microsoft Fabric} meters all workloads against a shared pool of Capacity Units
(CU) on an F-SKU, an approach closer to a single elastic budget than to per-engine warehouses.
\textbf{MongoDB Atlas} scales through cluster tiers with optional auto-scaling of compute and storage.

\section{Data Organization and the Medallion Pattern}
\index{medallion architecture!equivalence}%
Fabric's Chapter~2 topic---the Bronze/Silver/Gold medallion---is a portable design pattern rather
than a product. Every platform implements the same progressive-refinement idea with its own primitives.

\begin{center}
\footnotesize
\begin{tabular}{@{}p{2.6cm}p{3.4cm}p{3.4cm}p{3.0cm}@{}}
\toprule
\textbf{Layer} & \textbf{Fabric} & \textbf{Data lake (AWS/Azure/GCP)} & \textbf{Snowflake} \\
\midrule
Bronze (raw) & Bronze Lakehouse & \texttt{raw/} prefix (S3/ADLS/GCS) & RAW schema / stage \\
Silver (validated) & Silver Lakehouse & \texttt{validated/} prefix & STAGING schema \\
Gold (curated) & Gold Lakehouse & \texttt{curated/} prefix & ANALYTICS schema \\
Storage format & Delta-Parquet & Parquet / Delta / Iceberg & Micro-partitions \\
\bottomrule
\end{tabular}
\end{center}

\noindent \textbf{MongoDB} realizes the same separation with distinct collections or databases
(for example \texttt{raw\_events}, \texttt{validated}, \texttt{curated}) governed by schema-design
patterns rather than a lake layout.

\section{Document Data Modeling}
\index{document modeling!equivalence}%
MongoDB's Chapter~2 schema-design patterns generalize to every document and wide-column store.

\begin{center}
\footnotesize
\begin{tabular}{@{}p{3.0cm}p{3.2cm}p{3.2cm}p{3.0cm}@{}}
\toprule
\textbf{Concept} & \textbf{MongoDB} & \textbf{Amazon DynamoDB} & \textbf{Azure Cosmos DB} \\
\midrule
Primary access unit & Document / collection & Item / table & Item / container \\
Partitioning & Shard key & Partition key & Partition key \\
One-table modeling & Embedding + patterns & Single-table design & Embedding + partitioning \\
Time-series grouping & Bucket Pattern & Composite sort keys & Bucketing by key \\
Optional attributes & Attribute Pattern & Sparse attributes & Sparse properties \\
\bottomrule
\end{tabular}
\end{center}

\noindent Google Cloud's closest analogs are \textbf{Firestore} and \textbf{Bigtable}. In
\textbf{Snowflake} and \textbf{Fabric}, semi-structured documents are stored in \texttt{VARIANT}
(Snowflake) or JSON columns and queried with path expressions rather than modeled as collections.

\section{Globally Distributed and NoSQL Databases}
\index{NoSQL!equivalence}\index{distributed databases!equivalence}%
Chapter~3 branches by platform into high-scale NoSQL (AWS DynamoDB, Azure Cosmos DB), globally
consistent NewSQL (GCP Cloud Spanner), and advanced document modeling (MongoDB). They all answer
the same question: how to scale a database horizontally across partitions and regions.

\begin{center}
\footnotesize
\begin{tabular}{@{}p{2.8cm}p{2.9cm}p{2.9cm}p{2.9cm}p{2.4cm}@{}}
\toprule
\textbf{Concept} & \textbf{DynamoDB} & \textbf{Cosmos DB} & \textbf{Cloud Spanner} & \textbf{MongoDB} \\
\midrule
Model & Key-value / doc & Multi-model & Distributed SQL & Document \\
Partitioning & Partition key & Partition key & Split by key range & Shard key \\
Throughput unit & RCU / WCU & Request Units (RU) & Nodes / processing units & Cluster tier \\
Secondary indexes & GSI / LSI & Automatic indexing & Secondary indexes & Indexes \\
Change stream & DynamoDB Streams & Change Feed & Change streams & Change streams \\
Expiry & TTL & TTL & (manual) & TTL index \\
Multi-region & Global Tables & Multi-region writes & Multi-region (synchronous) & Zones / global clusters \\
Consistency & Eventual / strong & 5 tunable levels & External (TrueTime) & Tunable read/write concern \\
\bottomrule
\end{tabular}
\end{center}

\noindent \textbf{Snowflake} and \textbf{Fabric} are analytical platforms rather than OLTP key-value
stores; they ingest from these operational databases and store semi-structured payloads in
\texttt{VARIANT}/JSON columns. Google's dedicated NoSQL analogs are \textbf{Firestore} (document)
and \textbf{Bigtable} (wide-column).

\section{Managed Apache Spark}
\index{Apache Spark!equivalence}%
Microsoft Fabric's Chapter~3 topic---managed Spark pools and notebooks---has a first-class
equivalent on every cloud.

\begin{center}
\footnotesize
\begin{tabular}{@{}p{2.8cm}p{3.2cm}p{3.2cm}p{3.2cm}@{}}
\toprule
\textbf{Capability} & \textbf{Fabric} & \textbf{AWS} & \textbf{Azure / GCP} \\
\midrule
Managed Spark & Starter / Custom pools & Glue / EMR / EMR Serverless & Synapse Spark / Databricks; Dataproc \\
Notebook UX & Fabric Notebooks & Glue Studio / EMR Studio & Synapse / Databricks; Vertex AI \\
Serverless option & Autoscaling pools & Glue / EMR Serverless & Synapse serverless; Dataproc Serverless \\
Lake format & Delta on OneLake & Delta / Iceberg / Hudi & Delta; Delta / Iceberg \\
\bottomrule
\end{tabular}
\end{center}

\noindent \textbf{Snowflake} offers \textbf{Snowpark} for DataFrame-style processing, and
\textbf{MongoDB} integrates through the Spark Connector.

\section{Query Acceleration and Materialized Views}
\index{materialized views!equivalence}\index{query acceleration!equivalence}%
Snowflake's Chapter~3 acceleration tools---materialized views, the Search Optimization Service, and
transient/temporary tables---map to comparable features in every warehouse.

\begin{center}
\footnotesize
\begin{tabular}{@{}p{2.8cm}p{3.0cm}p{3.0cm}p{3.4cm}@{}}
\toprule
\textbf{Feature} & \textbf{Snowflake} & \textbf{AWS Redshift} & \textbf{Azure Synapse / GCP BigQuery} \\
\midrule
Materialized views & Materialized Views & Materialized views & Materialized views \\
Point-lookup accel. & Search Optimization Svc & Sort keys / zone maps & (indexing) / search indexes \\
Result reuse & Result cache & Result cache & Result-set cache / BI Engine \\
Ephemeral tables & Transient / temporary & Temp tables & Temp tables \\
\bottomrule
\end{tabular}
\end{center}

\noindent \textbf{Fabric} accelerates reads through Direct Lake and its Warehouse engine, while
\textbf{MongoDB} relies on secondary and compound indexes plus the in-memory working set.

\section{In-Memory Caching and Hybrid Ingestion}
\index{caching!equivalence}\index{in-memory cache!equivalence}%
The AWS and Azure Chapter~4 topic---low-latency caching in front of a durable store, fed by hybrid
(edge-to-cloud) ingestion---has a managed Redis-compatible equivalent on every cloud.

\begin{center}
\footnotesize
\begin{tabular}{@{}p{2.8cm}p{3.0cm}p{3.0cm}p{3.4cm}@{}}
\toprule
\textbf{Capability} & \textbf{AWS} & \textbf{Azure} & \textbf{GCP / others} \\
\midrule
Managed cache & ElastiCache (Redis/Memcached) & Cache for Redis & Memorystore (Redis/Memcached) \\
Real-time API layer & AppSync / API Gateway & API Management / SignalR & API Gateway / Apigee \\
Hybrid ingestion & DataSync / Snow family / IoT Core & Data Box / IoT Hub / Arc & Transfer Appliance / IoT Core \\
Streaming intake & Kinesis / MSK & Event Hubs & Pub/Sub \\
\bottomrule
\end{tabular}
\end{center}

\noindent \textbf{Snowflake} and \textbf{Fabric} rely on result and Direct~Lake caches rather than a
standalone Redis tier, while \textbf{MongoDB} serves hot data from its in-memory working set and
Atlas can pair with an external cache.

\section{Wide-Column and Key-Value NoSQL}
\index{wide-column!equivalence}\index{Bigtable!equivalence}%
GCP's Chapter~4 topic---Cloud Bigtable---is a wide-column, high-throughput, low-latency NoSQL store.
Its peers span the managed and open-source ecosystems.

\begin{center}
\footnotesize
\begin{tabular}{@{}p{2.8cm}p{3.0cm}p{3.0cm}p{3.4cm}@{}}
\toprule
\textbf{Characteristic} & \textbf{GCP Bigtable} & \textbf{AWS} & \textbf{Azure / open source} \\
\midrule
Data model & Wide-column & DynamoDB (key-value/doc) & Cosmos DB (multi-model) \\
HBase-compatible & HBase API & Keyspaces (Cassandra) & Managed Instance for Cassandra \\
Scaling unit & Nodes / clusters & On-demand / provisioned & RU/s (Cosmos) \\
Best fit & Time-series, IoT, ad-tech & Serverless key-value & Global multi-model \\
\bottomrule
\end{tabular}
\end{center}

\noindent \textbf{Snowflake} and \textbf{Fabric} address these workloads analytically rather than as
operational key-value stores; \textbf{MongoDB} covers the document end of the same NoSQL spectrum.

\section{Data Protection: Cloning, Time Travel, and Migration}
\index{cloning!equivalence}\index{time travel!equivalence}\index{point-in-time recovery!equivalence}%
Snowflake's Chapter~4 topic---zero-copy cloning, Time Travel, and Fail-safe---plus MongoDB's
relational-to-document migration map to data-protection and migration tooling across the clouds.

\begin{center}
\footnotesize
\begin{tabular}{@{}p{2.8cm}p{3.0cm}p{3.0cm}p{3.4cm}@{}}
\toprule
\textbf{Capability} & \textbf{Snowflake} & \textbf{AWS / Azure} & \textbf{GCP / Fabric} \\
\midrule
Instant clone & Zero-copy clone & Aurora clone / SQL DB copy & BigQuery table clone; OneLake shortcut \\
Point-in-time recovery & Time Travel & PITR (RDS/SQL DB backups) & PITR; Delta time travel \\
Retention safety net & Fail-safe (7 days) & Automated backups & Backups; Delta VACUUM window \\
Schema migration & --- & DMS / SCT & Database Migration Service \\
\bottomrule
\end{tabular}
\end{center}

\noindent \textbf{MongoDB} migrations use \textbf{Relational Migrator} plus the embed-vs-reference
patterns from Chapter~4, and \textbf{Delta Lake} exposes historical versions through
\texttt{VERSION AS OF} time-travel queries.

% ---------------------------------------------------------------------------
% Chapters 5--24: each section below harvests the cross-cloud callout of the
% named Snowflake chapter and adds the Databricks column from the companion
% Databricks book's cross-cloud boxes.
% ---------------------------------------------------------------------------

\section{Stages and Bulk File Ingestion}
\index{stage!equivalence}\index{bulk ingestion!equivalence}%
Snowflake's Chapter~5 topic---named, governed pointers to cloud object storage---rests on
primitives every provider offers: a bucket or container, an identity mechanism for a service to
read it, and an event channel that announces new files.

\begin{center}
\footnotesize
\begin{tabular}{@{}p{2.5cm}p{2.9cm}p{2.7cm}p{2.9cm}p{2.9cm}@{}}
\toprule
\textbf{Capability} & \textbf{Snowflake} & \textbf{AWS} & \textbf{Azure / GCP} & \textbf{Databricks} \\
\midrule
File location & External stage & S3 prefix & ADLS / GCS path & Volume / external location \\
Service credential & Storage integration & IAM role trust policy & Managed identity / service account & Storage credential + external location \\
Cloud-side grant & Integration-scoped locations & Bucket policy + role & RBAC / IAM binding & Unity Catalog grants on locations \\
New-file event & SQS / Event Grid / PubSub $\rightarrow$ Snowpipe & S3 event $\rightarrow$ SQS & Event Grid / Pub/Sub & OneLake events / file notifications \\
Bulk load & \texttt{COPY INTO} & Glue / Athena & ADF / BigQuery load & Auto Loader / \texttt{COPY INTO} \\
Offline transfer & Provider appliances & Snow Family & Data Box / Transfer Appliance & Same provider appliances \\
\bottomrule
\end{tabular}
\end{center}

\noindent \textbf{Microsoft Fabric} shortcuts virtualize ADLS Gen2, S3, and GCS inside OneLake
without copying---the same ``reference, do not duplicate'' idea as an external stage.
\textbf{MongoDB} has no object-stage concept; ingestion flows through drivers,
\texttt{mongoimport}, or Atlas Data Federation. The Databricks analog differs in one important
way: volumes and external locations are governed by Unity Catalog, so the pointer and its
permissions live in the same catalog as the tables the files become.

\section{Event-Driven and Streaming Ingestion}
\index{streaming ingestion!equivalence}\index{Snowpipe!equivalence}%
Snowflake's Chapter~6 pattern---an object lands, an event fires, a serverless loader reacts---is
universal; only the service names change.

\begin{center}
\footnotesize
\begin{tabular}{@{}p{2.5cm}p{2.9cm}p{2.7cm}p{2.9cm}p{2.9cm}@{}}
\toprule
\textbf{Capability} & \textbf{Snowflake} & \textbf{AWS} & \textbf{Azure / GCP} & \textbf{Databricks} \\
\midrule
Serverless file loader & Snowpipe auto-ingest & Lambda / Glue bookmarks & Function / Dataflow & Auto Loader (cloudFiles) \\
Managed streaming bus & (downstream consumer) & Kinesis / MSK & Event Hubs / Pub/Sub & Kafka connectors; Lakeflow \\
Row-streaming ingest & Snowpipe Streaming API & Firehose $\rightarrow$ Redshift & Event Hubs capture / Storage Write API & Structured Streaming sinks \\
Exactly-once mechanism & Channel offset tokens & Transactional sinks & Checkpointing / offsets & Checkpoint + write-ahead offsets \\
Malformed rows & \texttt{ON\_ERROR} modes & Dead-letter patterns & Error outputs / DLQ & \texttt{\_rescued\_data} column \\
Schema drift & File-format options / schema evolution & Glue schema detection & ADF mapping drift & Inference + hints + evolution \\
\bottomrule
\end{tabular}
\end{center}

\noindent \textbf{Microsoft Fabric} routes events through Eventstream into KQL databases or
lakehouse tables; \textbf{MongoDB} ingests through drivers and change streams rather than file
events. Snowflake's distinctive move is that the pipe is serverless---no cluster to size, billed
per second of pipe compute plus per-file notifications---while Databricks' Auto Loader trades
that serverless simplicity for checkpointed, schema-evolving streams on its own compute.

\section{Change Data Capture with Streams}
\index{change data capture!equivalence}\index{stream!equivalence}%
Snowflake's Chapter~7 idea---read the database's change journal once, replay it everywhere---is
one of the most portable concepts in data engineering.

\begin{center}
\footnotesize
\begin{tabular}{@{}p{2.5cm}p{2.9cm}p{2.7cm}p{2.9cm}p{2.9cm}@{}}
\toprule
\textbf{Capability} & \textbf{Snowflake} & \textbf{AWS} & \textbf{Azure / GCP} & \textbf{Databricks} \\
\midrule
Change journal object & Streams (standard / append-only / insert-only) & Kinesis Data Streams & Event Hubs / Pub/Sub & Delta change data feed \\
Managed CDC replication & (consumes via connectors) & AWS DMS & ADF CDC / Datastream & Lakeflow Connect, partner tools \\
Offset / checkpoint & Transactional stream offset & Kinesis sequence numbers & Checkpoints / subscriptions & Checkpointed stream offsets \\
CDC application & Stream + \texttt{MERGE} task & Glue/EMR \texttt{MERGE} & T-SQL / BigQuery \texttt{MERGE} & \texttt{MERGE} / DLT \texttt{AUTO CDC} \\
Staleness limit & 14-day stream retention & Stream retention windows & Retention per service & Log retention / \texttt{VACUUM} window \\
\bottomrule
\end{tabular}
\end{center}

\noindent \textbf{Microsoft Fabric} offers mirroring and event-driven pipelines; \textbf{MongoDB}
change streams on the oplog are the closest conceptual sibling to a Snowflake stream. Snowflake's
distinctive contribution is that the CDC journal is a queryable SQL object with a transactional
offset: consuming a stream and merging its deltas commit or roll back as one unit, whereas the
lakehouse family separates the change feed (Delta metadata) from the application step (a job).

\section{In-Database Scheduling and Task DAGs}
\index{task!equivalence}\index{orchestration!equivalence}%
Snowflake's Chapter~8 topic---scheduling and dependency graphs inside the database---is the
in-warehouse member of a family every cloud grew: a cron scheduler, a DAG engine, and a
serverless variant.

\begin{center}
\footnotesize
\begin{tabular}{@{}p{2.5cm}p{2.9cm}p{2.7cm}p{2.9cm}p{2.9cm}@{}}
\toprule
\textbf{Capability} & \textbf{Snowflake} & \textbf{AWS} & \textbf{Azure / GCP} & \textbf{Databricks} \\
\midrule
Cron scheduler & Tasks (\texttt{USING CRON}) & EventBridge Scheduler & Functions timer / Cloud Scheduler & Job schedules (cron) \\
DAG orchestrator & Task graphs via \texttt{AFTER} & Step Functions / MWAA & Data Factory / Cloud Composer & Workflows (jobs + tasks) \\
Conditional flow & \texttt{WHEN} predicates & Choice states & If Condition activities & if/else, for-each tasks \\
Serverless billing & Serverless tasks (per-second) & Lambda per invocation & Consumption plans / per-request & Serverless job compute \\
Partial recovery & Retry + \texttt{SUSPEND\_TASK\_AFTER\_NUM\_FAILURES} & Redrive / Airflow retries & Rerun from activity & Repair run (failed + downstream) \\
Failure notification & \texttt{ERROR\_INTEGRATION} & SNS & Azure Monitor / Pub/Sub & Job notifications / webhooks \\
\bottomrule
\end{tabular}
\end{center}

\noindent \textbf{Microsoft Fabric} orchestrates with Data Factory pipelines and notebook
scheduling; \textbf{MongoDB} relies on Atlas triggers and external schedulers. Snowflake's edge
is proximity: when the transformation is already SQL, orchestrating inside the database removes
an entire external system---its credentials, networking, and failure modes. Databricks'
Workflows keeps the orchestrator colocated with clusters and lineage, but it still schedules
external compute rather than warehouse-resident statements.

\section{Semi-Structured Data and Data Contracts}
\index{semi-structured data!equivalence}\index{data contract!equivalence}%
Snowflake's Chapter~9 answers three questions every platform faces: what type holds the document,
how do you project paths out of it, and where do you validate its shape.

\begin{center}
\footnotesize
\begin{tabular}{@{}p{2.5cm}p{2.9cm}p{2.7cm}p{2.9cm}p{2.9cm}@{}}
\toprule
\textbf{Capability} & \textbf{Snowflake} & \textbf{AWS} & \textbf{Azure / GCP} & \textbf{Databricks} \\
\midrule
Document column type & \texttt{VARIANT} / \texttt{OBJECT} / \texttt{ARRAY} & Redshift SUPER & OPENJSON / BigQuery JSON & \texttt{VARIANT} type; Delta JSON \\
Path extraction & Colon path notation & PartiQL dot/bracket & \texttt{JSON\_VALUE} / \texttt{JSON\_QUERY} & \texttt{get\_json\_object}, \texttt{:} operator \\
Array unnesting & \texttt{LATERAL FLATTEN} & PartiQL UNNEST & OPENJSON / UNNEST & \texttt{explode} / \texttt{inline} \\
Schema-on-read lake & External tables & Glue crawlers + Athena & Synapse serverless / external tables & Auto Loader inference; UC external tables \\
Contract enforcement & CHECK constraints, quarantine & Glue Schema Registry & Purview / Pub/Sub schemas & Delta constraints, DLT expectations \\
\bottomrule
\end{tabular}
\end{center}

\noindent \textbf{Microsoft Fabric} stores JSON in Delta/Parquet lakehouse files and explores it
with Spark or KQL; \textbf{MongoDB} \emph{is} the document model---its \texttt{\$jsonSchema}
validators are the closest analog to this chapter's contract checks. The shared lesson: one
flexible column or file type absorbs arbitrary payloads, while extracted attributes are cast,
constrained, and governed like relational columns.

\section{Query Profiling and Performance Tuning}
\index{query profiling!equivalence}\index{performance tuning!equivalence}%
Snowflake's Chapter~10 discipline---scan less, spill never, join carefully---is identical
everywhere; each platform exposes it through a different explain plan.

\begin{center}
\footnotesize
\begin{tabular}{@{}p{2.5cm}p{2.9cm}p{2.7cm}p{2.9cm}p{2.9cm}@{}}
\toprule
\textbf{Capability} & \textbf{Snowflake} & \textbf{AWS (Redshift)} & \textbf{Azure / GCP} & \textbf{Databricks} \\
\midrule
Plan inspection & Query Profile operators & EXPLAIN + monitoring views & DMVs / query plan + slot timeline & Spark UI + query profile \\
Data skipping & Micro-partition pruning metadata & Zone maps on sort keys & Partition pruning + clustering & File stats, ZORDER, liquid clustering \\
Spill signal & Local vs.\ remote spill bytes & Disk-based query steps & TempDB spill / bytes shuffled & Spill metrics in Spark UI \\
Adaptive execution & Automatic (engine-internal) & Auto WLM & Autoscale / managed & AQE on by default; Photon \\
Compute remedy & Warehouse resize & Resize / concurrency scaling & Scale DWUs / more slots & Cluster size, serverless, Photon \\
Result reuse & Persisted result cache & Result cache & Result-set / BI cache & Result cache, disk cache \\
\bottomrule
\end{tabular}
\end{center}

\noindent \textbf{MongoDB} answers the same questions through
\texttt{explain('executionStats')} and index design. Snowflake's twist is that there are no
indexes to build: the levers are pruning, warehouse size, caching, and SQL shape. Databricks
exposes the most knobs---AQE, join hints, Photon, shuffle partitions---and therefore the largest
tuning surface; BigQuery exposes the fewest, trading control for simplicity.

\section{Declarative Incremental Pipelines}
\index{dynamic tables!equivalence}\index{declarative pipelines!equivalence}%
Snowflake's Chapter~11 topic---declare freshness, let the platform schedule and refresh---is the
industry's answer to hand-built DAG sprawl.

\begin{center}
\footnotesize
\begin{tabular}{@{}p{2.5cm}p{2.9cm}p{2.7cm}p{2.9cm}p{2.9cm}@{}}
\toprule
\textbf{Capability} & \textbf{Snowflake} & \textbf{AWS} & \textbf{Azure / GCP} & \textbf{Databricks} \\
\midrule
Declarative refresh object & Dynamic tables (\texttt{TARGET\_LAG}) & Redshift MVs (auto-refresh) & Fabric pipelines / BigQuery MVs + scheduled queries & Lakeflow Declarative Pipelines (DLT) \\
Streaming dataset & (via dynamic table refresh cadence) & Redshift streaming MVs & Fabric Real-Time / continuous queries & DLT streaming tables \\
Quality rules & DMFs, tests downstream & Glue Data Quality (Deequ) & Great Expectations / dbt tests & DLT expectations (warn/drop/fail) \\
DAG derivation & Automatic from query text & Glue workflows (explicit) & Dataform dependency graph & Automatic from dataset definitions \\
Ops telemetry & \texttt{DYNAMIC\_TABLE\_REFRESH\_HISTORY} & CloudWatch + job metrics & ADF / Dataflow monitoring & DLT event log (Delta table) \\
\bottomrule
\end{tabular}
\end{center}

\noindent \textbf{MongoDB} offers on-demand materialized views via \texttt{\$merge}. dbt users
recognize the philosophy immediately: a dynamic table or DLT dataset is a dbt model whose
scheduler, dependency graph, and incremental state live inside the platform. The cross-cloud
lesson is that \emph{declared freshness} (target lag) is replacing \emph{declared schedule}
(cron) as the pipeline contract.

\section{Window Functions and Dimensional Modeling}
\index{window functions!equivalence}\index{dimensional modeling!equivalence}%
Snowflake's Chapter~12 skills are the most portable in the book: ISO SQL window semantics and the
Kimball method travel unchanged; only convenience clauses differ.

\begin{center}
\footnotesize
\begin{tabular}{@{}p{2.5cm}p{2.9cm}p{2.7cm}p{2.9cm}p{2.9cm}@{}}
\toprule
\textbf{Capability} & \textbf{Snowflake} & \textbf{AWS (Redshift)} & \textbf{Azure / GCP} & \textbf{Databricks} \\
\midrule
Window functions & Full ISO support & Full ISO support & Full ISO support & Full ISO support (Spark SQL) \\
Filter window output & \texttt{QUALIFY} & Subquery wrapper & Subquery / BigQuery \texttt{QUALIFY} & Subquery wrapper \\
Multi-grain aggregation & GROUPING SETS / CUBE / ROLLUP & Same & Same & Same \\
Pivot & PIVOT / UNPIVOT & CASE / PIVOT variants & PIVOT operators & PIVOT (Spark SQL) \\
Dimensional physical design & Pruning on fact filters & Sort/dist keys & Distribution + partition / partition + cluster & ZORDER / liquid clustering on dims \\
Slowly changing dims & \texttt{MERGE}-driven Type 2 & \texttt{MERGE} / staging & \texttt{MERGE} variants & \texttt{MERGE INTO} + change feed \\
\bottomrule
\end{tabular}
\end{center}

\noindent \textbf{MongoDB} expresses the same analytics in aggregation pipelines
(\texttt{\$group}, \texttt{\$setWindowFields}). The durable insight: star schemas remain the
serving layer of choice because denormalized dimensions match how columnar engines---warehouse
or lakehouse alike---prune and join.

\section{DataFrame APIs and Pushdown Execution}
\index{DataFrame API!equivalence}\index{Snowpark!equivalence}%
Snowflake's Chapter~13 topic---a client DataFrame API that compiles to the engine's native
execution---exists on every cloud; the API is the portable skill, the pushdown target is the
vendor's edge.

\begin{center}
\footnotesize
\begin{tabular}{@{}p{2.5cm}p{2.9cm}p{2.7cm}p{2.9cm}p{2.9cm}@{}}
\toprule
\textbf{Capability} & \textbf{Snowflake} & \textbf{AWS} & \textbf{Azure / GCP} & \textbf{Databricks} \\
\midrule
DataFrame API & Snowpark (Python/Scala/Java) & Spark on EMR / Glue & Spark on Synapse/Fabric; BigFrames & Apache Spark (reference implementation) \\
Pushdown target & Virtual warehouse (SQL) & Redshift Spectrum / Athena & Fabric SQL endpoint / BigQuery SQL & Photon / Spark on Delta \\
Lazy plan inspection & \texttt{df.explain()} & Spark explain & Spark explain & Spark explain / query profile \\
Pandas interop & Snowpark pandas, \texttt{to\_pandas} & pandas on SageMaker & pandas in notebooks & pandas API on Spark \\
Cluster management & None (warehouse size only) & EMR/Glue cluster sizing & Pool sizing & Cluster policies, autoscaling, serverless \\
\bottomrule
\end{tabular}
\end{center}

\noindent \textbf{MongoDB} exposes aggregation pipelines through language drivers---a different
algebra, the same client-side-builder pattern. Snowpark's defining property is that the data
never leaves Snowflake: the client builds an expression tree and the warehouse executes it as
SQL beside the data. Databricks is the reference implementation of the Spark API the others
host; Snowpark is the same client-side philosophy pointed at a warehouse instead of a cluster.

\section{User-Defined Functions and In-Engine Code}
\index{UDF!equivalence}\index{in-database code!equivalence}%
Snowflake's Chapter~14 topic---running user code where the data lives---is the great warehouse
convergence of recent years.

\begin{center}
\footnotesize
\begin{tabular}{@{}p{2.5cm}p{2.9cm}p{2.7cm}p{2.9cm}p{2.9cm}@{}}
\toprule
\textbf{Capability} & \textbf{Snowflake} & \textbf{AWS (Redshift)} & \textbf{Azure / GCP} & \textbf{Databricks} \\
\midrule
Scalar UDF & SQL / Python UDFs & SQL / Python UDFs & SQL-CLR lineage / BigQuery SQL \& JS UDFs & Spark UDFs (SQL/Python/Scala) \\
Vectorized UDF & Vectorized pandas UDFs (Arrow) & Lambda UDFs & Spark pandas UDFs / remote functions & pandas UDFs over Arrow \\
Table function & UDTFs & SQL UDFs returning tables & OPENJSON-style TVFs / table functions & Table-valued functions, generators \\
Package ecosystem & Curated Anaconda channel & Limited library imports & Spark environment packages / containers & Full PyPI on clusters; serverless environments \\
Runtime sandbox & Per-warehouse sandboxed runtime & Lambda sandbox & Spark executor / remote container & Spark executors / serverless \\
\bottomrule
\end{tabular}
\end{center}

\noindent \textbf{MongoDB} runs server-side JavaScript functions (legacy) and Atlas triggers.
Snowflake's differentiator is the sandboxed Python runtime with a curated Anaconda channel
inside the warehouse---no cluster to manage, and RBAC on the function itself. Databricks offers
the full PyPI ecosystem but makes you own the cluster the code runs on.

\section{Stored Procedures and In-Warehouse ML}
\index{stored procedure!equivalence}\index{ML in-database!equivalence}%
Snowflake's Chapter~15 modernizes the oldest in-database programmability---procedures---and its
ML variant trains models where the features already live.

\begin{center}
\footnotesize
\begin{tabular}{@{}p{2.5cm}p{2.9cm}p{2.7cm}p{2.9cm}p{2.9cm}@{}}
\toprule
\textbf{Capability} & \textbf{Snowflake} & \textbf{AWS (Redshift)} & \textbf{Azure / GCP} & \textbf{Databricks} \\
\midrule
Procedure language & SQL / Python procedures & PL/pgSQL & T-SQL / BigQuery scripting & Notebooks, SQL scripts, jobs \\
ML training in-engine & Snowpark ML (scikit-learn et al.) & Redshift ML (SageMaker) & Fabric notebooks / BigQuery ML & Full ML stack on Spark clusters \\
Experiment tracking & (external / registry hooks) & SageMaker Experiments & Azure ML / Vertex & Managed MLflow \\
Model registry & Snowflake Model Registry & SageMaker Model Registry & Azure ML / Vertex registries & Unity Catalog model registry \\
Feature store & Snowflake Feature Store & SageMaker Feature Store & Azure ML / Vertex feature stores & Feature Engineering in UC \\
\bottomrule
\end{tabular}
\end{center}

\noindent \textbf{MongoDB} runs server-side logic via Atlas triggers and functions. Snowflake's
proposition for ML teams: scikit-learn trains where the features already are, governed by the
same RBAC as the tables, with zero export pipeline to secure. Databricks answers with the
broader MLOps surface---MLflow, feature store, serving---at the cost of managing compute;
Snowflake answers with governed simplicity inside the warehouse.

\section{External Functions and Reverse ETL}
\index{external function!equivalence}\index{reverse ETL!equivalence}%
Snowflake's Chapter~16 integration pair---calling HTTP services from inside the warehouse, and
pushing warehouse results back to operational tools---exists in both directions on every cloud.

\begin{center}
\footnotesize
\begin{tabular}{@{}p{2.5cm}p{2.9cm}p{2.7cm}p{2.9cm}p{2.9cm}@{}}
\toprule
\textbf{Capability} & \textbf{Snowflake} & \textbf{AWS} & \textbf{Azure} & \textbf{GCP / Databricks} \\
\midrule
API front door & API integration & API Gateway & API Management & API Gateway / Endpoints \\
Compute behind API & External function $\rightarrow$ remote service & Lambda & Azure Functions & Cloud Run / Functions \\
Warehouse-to-API call & External functions & Redshift Lambda UDFs & Synapse + Logic Apps & BigQuery remote functions; Databricks SQL UDFs calling endpoints \\
Reverse ETL & Task + stream + external function & Census/Hightouch + EventBridge & Census/Hightouch + ADF & Census/Hightouch + Functions; Databricks jobs + webhooks \\
Third-party auth & Secrets + network rules & IAM + secrets in API layer & Managed identity + Key Vault & Service accounts + Secret Manager; Databricks secret scopes \\
\bottomrule
\end{tabular}
\end{center}

\noindent \textbf{MongoDB} reaches outward via Atlas triggers and HTTPS endpoints;
\textbf{Fabric} via notebooks and pipelines. The Snowflake constraint to remember: external
functions are synchronous and billed per call---the architecture must batch and cache, or the
API bill becomes the story. The same cost shape applies anywhere a query engine synchronously
invokes a metered endpoint.

\section{RBAC, Networking, and Secrets}
\index{RBAC!equivalence}\index{network security!equivalence}\index{secrets management!equivalence}%
Snowflake's Chapter~17 concentrates identity, role hierarchies, and network perimeters inside
the data platform; the disciplines themselves are cloud-native.

\begin{center}
\footnotesize
\begin{tabular}{@{}p{2.5cm}p{2.9cm}p{2.7cm}p{2.9cm}p{2.9cm}@{}}
\toprule
\textbf{Capability} & \textbf{Snowflake} & \textbf{AWS} & \textbf{Azure / GCP} & \textbf{Databricks} \\
\midrule
Workforce identity & SSO/SCIM, MFA, key-pair auth & IAM Identity Center & Entra ID / Cloud Identity & SSO + SCIM from IdP \\
Privilege grouping & Role hierarchy (access vs.\ functional) & IAM policies \& roles & RBAC / IAM roles \& groups & UC grants + account groups \\
Machine identity & Service users + key pairs & IAM roles & Managed identity / service accounts & Service principal + OAuth M2M \\
Long-lived secret & (discouraged; key rotation) & Access keys & Client secrets / SA keys & PAT (policy-controlled) \\
Network perimeter & Network policies, private connectivity & VPC endpoints, SCPs & Private Link / VPC Service Controls & Private connectivity, IP access lists \\
Secrets store & External (cloud KMS/vault) & Secrets Manager & Key Vault / Secret Manager & Secret scopes (DB- or vault-backed) \\
Audit & \texttt{ACCOUNT\_USAGE} views & CloudTrail & Activity Log / Cloud Audit Logs & \texttt{system.access.audit} \\
\bottomrule
\end{tabular}
\end{center}

\noindent \textbf{Microsoft Fabric} layers workspace roles over OneLake security;
\textbf{MongoDB} uses database roles and Atlas IP access lists. Snowflake's distinctive rule:
privileges are granted to roles, roles to roles, and roles to users---never privileges directly
to users---so access review becomes a graph query. Unity Catalog's distinctive move is one
governance model spanning tables, files, notebooks, models, and functions.

\section{Masking, Row Access, and Classification Tags}
\index{masking policy!equivalence}\index{row access policy!equivalence}\index{data classification!equivalence}%
Snowflake's Chapter~18 pattern---tag the data once, let policies follow the tags---is the
industry's classification-driven enforcement model.

\begin{center}
\footnotesize
\begin{tabular}{@{}p{2.5cm}p{2.9cm}p{2.7cm}p{2.9cm}p{2.9cm}@{}}
\toprule
\textbf{Capability} & \textbf{Snowflake} & \textbf{AWS} & \textbf{Azure / GCP} & \textbf{Databricks (UC)} \\
\midrule
Classification & Object tags + discovery & Macie + resource tags & Purview labels / Sensitive Data Protection & Tags + PII discovery \\
Column masking & Masking policies (tag-based) & Redshift dynamic masking / Lake Formation & Fabric masking / BigQuery policy tags & Column masks \\
Row-level security & Row access policies & Lake Formation row filters & Synapse RLS / BigQuery row access & Row filters \\
Policy-tag binding & Tag $\rightarrow$ policy once & LF tag expressions & Label-inherited protection / policy taxonomy & Tag-driven ABAC policies \\
Erasure support & Crypto-shredding (per-subject keys) & KMS shred / Object Lock exceptions & CMK shredding / CMEK destruction & CMK + GDPR deletes on Delta \\
\bottomrule
\end{tabular}
\end{center}

\noindent \textbf{MongoDB} offers client-side field-level encryption; \textbf{Fabric} inherits
Purview labels into OneLake. Snowflake's distinctive leverage: tags attach to columns and
policies attach to tags, so classifying a column once protects it in every query, clone, and
share that touches it. Databricks reaches the same outcome through Unity Catalog attribute-based
access control over its catalogs.

\section{Secure Data Sharing and Marketplaces}
\index{data sharing!equivalence}\index{marketplace!equivalence}%
Snowflake's Chapter~19 idea---sharing analytics data without copying it---is warehouse-native;
other platforms approximate it with exports, replicas, or open sharing protocols.

\begin{center}
\footnotesize
\begin{tabular}{@{}p{2.5cm}p{2.9cm}p{2.7cm}p{2.9cm}p{2.9cm}@{}}
\toprule
\textbf{Capability} & \textbf{Snowflake} & \textbf{AWS} & \textbf{Azure / GCP} & \textbf{Databricks} \\
\midrule
In-platform sharing & Secure shares (metadata only) & Redshift data sharing & Fabric shortcuts / BigQuery Analytics Hub & Delta Sharing (open protocol) \\
Non-customer access & Reader accounts & S3 presigned / Spectrum & Data Share invitations / authorized views & Any recipient (pandas, Spark, Power BI) \\
Data marketplace & Snowflake Marketplace & AWS Data Exchange & Azure Marketplace / Analytics Hub & Databricks Marketplace \\
Consumer pays compute & Consumer account or reader account & Redshift consumer cluster & Fabric capacity / consumer slots & Recipient's own compute \\
Shared payload & Secure views, tables, UDFs & Datashares & Shortcuts / listings & Tables, volumes, notebooks, models \\
\bottomrule
\end{tabular}
\end{center}

\noindent \textbf{MongoDB} shares via Atlas exports and API layers. Snowflake's distinction is
that sharing is \emph{metadata}: the consumer queries the provider's micro-partitions live with
the provider's governance attached. Delta Sharing's distinction is the \emph{open} protocol:
recipients consume without any Databricks account. Both replace the export file; they differ on
who must run whose platform.

\section{Data Clean Rooms and Native Applications}
\index{clean room!equivalence}\index{native app!equivalence}%
Snowflake's Chapter~20 topics---privacy-preserving collaboration and shipping applications to
consumers---have been productized by every hyperscaler with different seams.

\begin{center}
\footnotesize
\begin{tabular}{@{}p{2.5cm}p{2.9cm}p{2.7cm}p{2.9cm}p{2.9cm}@{}}
\toprule
\textbf{Capability} & \textbf{Snowflake} & \textbf{AWS} & \textbf{Azure / GCP} & \textbf{Databricks} \\
\midrule
Clean room service & Snowflake Data Clean Rooms & AWS Clean Rooms & Confidential computing / BigQuery clean rooms & Delta Sharing clean rooms \\
Identifier protection & Hashed/keyed join keys & Hashed/normalized keys & Enclaves / confidential space & Hashed keys over shared tables \\
Aggregate thresholds & Minimum cohort rules & Minimum cohort rules & Policy / analysis rules & Policy-enforced in clean room \\
Differential privacy & Clean-room DP options & AWS Clean Rooms DP & SmartNoise / BigQuery DP & (partner / library-based) \\
App distribution & Native App Framework & Marketplace AMIs/SaaS & Azure / GCP Marketplaces & Marketplace + apps (limited) \\
\bottomrule
\end{tabular}
\end{center}

\noindent \textbf{MongoDB} and \textbf{Fabric} rely on partner architectures rather than native
clean-room products. Snowflake's compounding advantage: clean rooms build on zero-copy sharing,
so multi-party analysis never moves anyone's data, and Native Apps ship not just data but the
code that consumes it---installed into the consumer's account and run on their compute.

\section{HTAP and Hybrid Tables}
\index{HTAP!equivalence}\index{hybrid tables!equivalence}%
Snowflake's Chapter~21 question---can one store serve both transactional and analytical
workloads?---has an answer on every cloud.

\begin{center}
\footnotesize
\begin{tabular}{@{}p{2.5cm}p{2.9cm}p{2.7cm}p{2.9cm}p{2.9cm}@{}}
\toprule
\textbf{Capability} & \textbf{Snowflake} & \textbf{AWS} & \textbf{Azure / GCP} & \textbf{Databricks} \\
\midrule
HTAP offering & Hybrid tables (Unistore) & Aurora $\rightarrow$ Redshift zero-ETL & Cosmos + Synapse Link / AlloyDB columnar & Lakebase (managed OLTP on Postgres) \\
Row-store + analytics & One platform, both engines & OLTP + zero-ETL replica & Analytical store replicas & Delta for analytics; Lakebase for OLTP \\
Constraint enforcement & Enforced PK/FK/unique on hybrid tables & Full relational DDL & Full relational DDL & Delta constraints (limited); Lakebase full \\
Point-read latency & Millisecond (row store) & Millisecond (OLTP engine) & Millisecond & Millisecond (Lakebase) / seconds (Delta) \\
Serving layer & Hybrid tables + same SQL & Aurora/RDS serving & Cosmos/Spanner serving & Online tables / feature serving \\
\bottomrule
\end{tabular}
\end{center}

\noindent \textbf{MongoDB} is the original operational/analytics hybrid discussion (replica-set
analytics nodes); \textbf{Fabric} separates KQL/warehouse engines from operational stores.
Snowflake's Unistore move: the transactional table lives inside the analytical database, so
joining orders to the warehouse mart is a same-platform query, not a replication pipeline.
Databricks answers with a managed Postgres (Lakebase) beside the lakehouse rather than a second
table type inside it.

\section{Managed AI Functions, Vectors, and Document Intelligence}
\index{Cortex!equivalence}\index{vector search!equivalence}\index{LLM functions!equivalence}%
Snowflake's Chapter~22 topic---LLM and embedding functions inside the data platform---is now
table stakes; the differences are model catalogs, billing shape, and how natively vectors live
in the engine.

\begin{center}
\footnotesize
\begin{tabular}{@{}p{2.5cm}p{2.9cm}p{2.7cm}p{2.9cm}p{2.9cm}@{}}
\toprule
\textbf{Capability} & \textbf{Snowflake} & \textbf{AWS} & \textbf{Azure / GCP} & \textbf{Databricks} \\
\midrule
Managed LLM service & \texttt{CORTEX.*} functions & Bedrock & Azure OpenAI / Vertex AI & Foundation Model APIs \\
In-database LLM call & SQL functions on columns & Redshift ML $\rightarrow$ Bedrock & Notebooks + AOAI / BigQuery ML remote models & \texttt{ai\_query} in SQL/Spark \\
Embeddings & \texttt{EMBED\_TEXT\_768} & Bedrock Titan embeddings & AOAI / Vertex embeddings & Hosted embedding endpoints \\
Vector search & \texttt{VECTOR} type + similarity functions & OpenSearch / pgvector & AI Search / Vertex Vector Search & Mosaic AI Vector Search \\
Document parsing & \texttt{PARSE\_DOCUMENT} & Textract & Document Intelligence / Document AI & \texttt{ai\_parse\_document} \\
RAG orchestration & Cortex Search / custom SQL & Bedrock Knowledge Bases & Azure AI RAG / Vertex RAG Engine & LangChain + MLflow, Agent Bricks \\
\bottomrule
\end{tabular}
\end{center}

\noindent \textbf{MongoDB} Atlas Vector Search pairs operational documents with embeddings;
\textbf{Fabric} routes through Azure OpenAI. Snowflake's distinction is governance continuity:
documents, embeddings, prompts, and answers live under the same RBAC and masking as every other
table---no extracts to a side vector database. Databricks' distinction is openness: any model,
any framework, governed by Unity Catalog over chunks and endpoints.

\section{DataOps: IaC, Testing, and Observability}
\index{DataOps!equivalence}\index{infrastructure as code!equivalence}\index{observability!equivalence}%
Snowflake's Chapter~23 is the most cloud-portable chapter in the book: Terraform, dbt, and
GitHub Actions are the same tools on every platform; only the provider plugin and the
warehouse-specific monitors differ.

\begin{center}
\footnotesize
\begin{tabular}{@{}p{2.5cm}p{2.9cm}p{2.7cm}p{2.9cm}p{2.9cm}@{}}
\toprule
\textbf{Capability} & \textbf{Snowflake} & \textbf{AWS} & \textbf{Azure / GCP} & \textbf{Databricks} \\
\midrule
Cross-platform IaC & terraform-provider-snowflake & Terraform AWS provider & azurerm / google providers & Terraform databricks provider \\
Native IaC & (SQL DDL, git-managed) & CloudFormation / CDK & Bicep / Deployment Manager & Asset Bundles (\texttt{databricks.yml}) \\
Transformation + tests & dbt-snowflake & dbt-redshift & dbt-fabric / dbt-bigquery & dbt-databricks / dbt-spark \\
Migration runner & schemachange / Flyway & Flyway / Liquibase & Flyway / DbUp & Asset Bundles / notebooks in CI \\
Quality monitors & DMFs, alerts, resource monitors & Glue DQ + CloudWatch & Purview / Dataplex monitors & Lakehouse Monitoring \\
Spend telemetry & \texttt{ACCOUNT\_USAGE}, monitors & CUR 2.0 + Cost Explorer & Cost Management / billing export & \texttt{system.billing.usage} \\
CI authentication & Key-pair secrets & OIDC $\rightarrow$ IAM role & OIDC $\rightarrow$ workload identity & Service principal + OAuth M2M \\
\bottomrule
\end{tabular}
\end{center}

\noindent \textbf{MongoDB} uses the same Terraform discipline with Atlas providers;
\textbf{Fabric} adds workspace deployment pipelines. The lesson that transfers everywhere:
\emph{the merge is the deploy, and the pipeline is the only writer to production}. The
Snowflake-specific craft is platform-native observability---resource monitors, alerts, DMFs, and
\texttt{ACCOUNT\_USAGE} as the telemetry source of truth.

\section{Certification and Career Capstone}
\index{certification!equivalence}%
Every platform book in this series ends at a professional certification, and the preparation
method---domains mapped to labs, timed mocks, error-driven remediation---is identical across
vendors.

\begin{center}
\footnotesize
\begin{tabular}{@{}p{2.5cm}p{2.9cm}p{2.7cm}p{2.9cm}p{2.9cm}@{}}
\toprule
\textbf{Dimension} & \textbf{Snowflake} & \textbf{AWS} & \textbf{Azure / GCP} & \textbf{Databricks} \\
\midrule
Entry certification & SnowPro Core & Cloud Practitioner / Data Engineer Assoc. & DP-900 / Cloud Digital Leader & Data Engineer Associate \\
Advanced certification & SnowPro Advanced: Data Engineer & Data Engineer Associate & DP-203/DP-600 / Professional Data Engineer & Data Engineer Professional \\
Exam style & Scenario MCQ & Scenario MCQ & Case studies / scenario MCQ & Scenario MCQ + hands-on fluency \\
Heaviest domains & Data movement, performance & Ingestion, storage, processing & Design + implementation / processing & Delta Lake, Spark ETL, pipelines \\
Mock standard & 80\%+, no domain weak & Same & Same & 80\%+, no domain $<$70\% \\
Hands-on evidence & Portfolio of capstones & Portfolio projects & Portfolio projects & Architecture review with cost numbers \\
\bottomrule
\end{tabular}
\end{center}

\noindent The durable artifact is never the badge: it is the portfolio of measured labs,
architecture diagrams, and cost models built across the program. If you can defend the
Snowflake capstone \emph{and} sketch its Databricks or AWS equivalent from the tables in this
appendix, you have learned data engineering, not one console.

\begin{tipbox}
Use this appendix as a study aid for the book's lockstep, cross-cloud approach: when you learn a
concept on one platform, immediately locate its equivalent here and predict how the other six
would express the same idea. That habit turns seven separate vocabularies into one mental model.
\end{tipbox}
